\documentclass[11pt,a4paper]{article}
\usepackage[margin=2.5cm]{geometry}
\usepackage[T1]{fontenc}
\usepackage{lmodern}
\usepackage{amsmath,amssymb}
\usepackage{booktabs}
\usepackage[hidelinks]{hyperref}
\usepackage{xcolor}
\newcommand{\code}[1]{\texttt{#1}}
\setlength{\parskip}{0.5em}
\setlength{\parindent}{0pt}

\title{Tracer against Medis: the comparison step by step}
\author{Working note, \code{CGPS\_data/thesis\_final/medis\_compare.py}}
\date{Code state of 19 November 2026 (after the code review of that day; not yet run on the server)}

\begin{document}
\maketitle

\noindent This note gives every step the script takes, in the order it takes them, with the constants
it uses. It describes the code, not the thesis text. Section~\ref{sec:changes} lists what changed on
19 November 2026. Those changes follow the code review and were made after the 18 November results had
been seen.

\begin{table}[h]
\centering
\begin{tabular}{@{}llp{8.2cm}@{}}
\toprule
Constant & Value & Meaning \\
\midrule
\code{MIN\_SHARED} & 10 mm & minimum shared stretch, and minimum averaged part \\
\code{CHORD} & 5 mm & chord of the curvature profile; also cut from each end of the shared stretch \\
\code{DENSE} & 0.25 mm & resampling step for the distances \\
\code{MAX\_OFFSET} & 2 mm & alignment cut on the 90th percentile of the distances \\
\code{DEPTHS} & 5, 5.5, \dots, 10 mm & depths from O used for the angle \\
\code{O\_FAR} & 5 mm & threshold for reporting $o$ (not a cut) \\
bootstrap & 2000 samples, seed 0 & resampling of participants \\
\bottomrule
\end{tabular}
\end{table}

\section{Input}

\textbf{Step 1: participants.} The input is all Medis centerline trees (\code{Medis/.../*.vtk}, 600 at
baseline), sorted by study id. A participant goes on only when both of the tracer's pruned baseline files
exist (\code{BL\_<id>.coronary\_\{left,right\}\_centerline\_ostium\_pruned.vtk}); otherwise the flow
counts the participant as ``no tracer pair''.

\textbf{Step 2: reading the trees.} The Medis tree is read as polylines with a segment label per point
(\code{hoe10y\_qc.tree\_points}; the label codes are named by \code{read\_label\_codes}). The tracer tree
is read as one polyline per branch of both files, with the voted integer code of each point mapped to the
same label names (\code{tracer\_io.tracer\_tree}, map \code{tracer\_io.CODES}). Both are taken as LPS mm
in the same frame, and no mapping is applied. (Which coordinates the two store is still open.) The
tracer's ostia are also read: the first point of the root branch (\code{parent} $=-1$) of the left file is
the left ostium, used for the LAD and the LCx, and that of the right file is the RCA ostium
(\code{tracer\_io.tracer\_ostia}).

\section{Dividing each tree into arteries}

The same function, \code{split\_trunks.split}, is applied to both trees. For the comparison only the
whole artery is used, not the segment boundaries the function also returns.

\textbf{Step 3: graph.} All points are pooled. Points that agree to 0.01 mm become one node, and a node
takes the first label any of its points has. Consecutive points of a polyline are joined by an edge.
A polyline end within 1 mm of a point of another polyline is joined to it. After that, a polyline end
within 3 mm (\code{BRIDGE}) of another connected piece is joined to the nearest node of that piece. The
edge weight is the Euclidean length (\code{split\_trunks.graph}).

\textbf{Step 4: path of each artery.} For the LAD, the LCx and the RCA in turn, the nodes labelled with
that artery (any of its segment labels) are collected. The connected piece of the graph that holds most
of them is the one used.

The root of the path is found differently for the two trees:
\begin{itemize}
\item \emph{Tracer:} the node nearest the tracer's own ostium from Step 2.
\item \emph{Medis:} the most superior (largest $z$) graph end labelled \code{lm}. Without one, it is the
most superior end labelled as the artery's proximal segment, then an artery or unlabelled end, then
any end.
\end{itemize}

The path is the shortest path (Dijkstra) from the root to the farthest node, in path length, that
carries the artery's distal label (\code{dlad}; \code{dcx} or \code{cx}; \code{drca}). Without such a
node, it runs to the farthest node labelled with the artery (\code{split\_trunks.trunk\_path}).

\textbf{Step 5: start at the bifurcation point O.} The LAD and LCx paths both start at the left ostium,
so they share their first nodes. When they share a start, both are cut at their last shared node. This
node is the bifurcation point O, and both daughters start exactly there. Only when they share no start
(for example one of them is missing) does a path start at its last node labelled \code{lm}
(\code{split\_trunks.start\_at\_bifurcation}). The RCA starts at its ostium. Each path gets the arc length
$s$ along it, 0 at its start.

A path with fewer than three nodes, or one that raises an error, counts as missing for that artery. A
tree on which the whole step raises is counted as ``tree not divided''.

\section{Curvature, per artery pair}

\textbf{Step 6: pairing.} The tracer LAD is paired with the Medis LAD of the same participant, and
likewise for the LCx and the RCA. An artery missing in one or both trees is kept as a row with the reason
``missing in Medis / tracer''. Every following step runs inside a guard: a pair whose computation raises
is kept with the reason \code{error:<type>}, and the run goes on.

\textbf{Step 7: shared stretch.} Let the tracer line be $Y_t$ with arc length $s_t\in[0,L_t]$, and the
Medis line $Y_m$ with $s_m\in[0,L_m]$. Write $\pi_t(q)$ for the arc length of the point on the tracer
polyline closest to $q$, measured on the line pieces and not only on the points; $\pi_m$ is the same on
Medis. The shared stretch is
\begin{align*}
[t_0,t_1]&=\bigl[\max(0,\pi_t(Y_m(0))),\ \min(L_t,\pi_t(Y_m(L_m)))\bigr],\\
[m_0,m_1]&=\bigl[\max(0,\pi_m(Y_t(0))),\ \min(L_m,\pi_m(Y_t(L_t)))\bigr],
\end{align*}
that is, from the later start to the earlier end, on each line (\code{shared\_stretch}). If
$\min(t_1-t_0,\ m_1-m_0)<10$ mm, the pair is left out with the reason ``shared stretch $<$ 10 mm''. It is
then not aligned.

\textbf{Step 8: alignment.} The two shared stretches are each resampled every 0.25 mm along their arc
length, by linear interpolation with both ends kept (\code{stretch}). This gives $A$ (tracer) and $B$
(Medis). For each point of $A$, the distance to the nearest point of $B$ is taken, and the same is done
from $B$ to $A$. The offset is
\[
\text{offset}=\max\bigl(P_{90}\{d(a,B)\}_{a\in A},\ P_{90}\{d(b,A)\}_{b\in B}\bigr).
\]
The pair is aligned (\code{aligned\_artery}) when the offset is at most 2 mm. The median of $d(a,B)$
(tracer to Medis) is saved as \code{offset\_median}, for the report only. Both directions are measured
only within the shared stretches, so an end cannot match a part of the other line outside the overlap.
Both are evenly resampled, so no stretch counts more because it has more points.

\textbf{Step 9: averaged part.} The averaged part is the shared stretch without 5 mm (\code{CHORD}) at
each end, on both lines: $[t_0+5,\ t_1-5]$ and $[m_0+5,\ m_1-5]$. If either is under 10 mm (a shared
stretch under 20 mm), the pair is left out with the reason ``averaged part $<$ 10 mm''. This reason comes
before the offset, so such a pair gets no curvature even if it also failed Step 8. Its
\code{aligned\_artery} from Step 8 is kept, because the angle uses it (Step 13).

\textbf{Step 10: reason.} A pair that reaches this step gets the reason ``ok'' when it is aligned and
``offset $>$ 2 mm'' otherwise. Curvature is computed for both kinds.

\textbf{Step 11: curvature.} A curvature profile is computed on each \emph{whole} artery, from O (or the
RCA ostium) to its end, not on the shared stretch alone. At each point $P_i$, the chord ends $B_i$ and
$F_i$ are the points on the polyline exactly 5 mm away from $P_i$ in a straight line, one behind and one
ahead, found by linear interpolation on the first piece that crosses 5 mm. The curvature is Menger's
\[
\kappa_i=\frac{2\,\lVert (P_i-B_i)\times(F_i-P_i)\rVert}{\lVert P_i-B_i\rVert\,\lVert F_i-P_i\rVert\,\lVert F_i-B_i\rVert}
\quad[\text{mm}^{-1}],
\]
the inverse radius of the circle through the three points. Points without both chord ends get no
profile value (\code{curvature\_feature.curvature\_profile}). The feature of each line is the arc-length
mean, by the trapezoid rule, of its profile values at $s\in[\text{lo},\text{hi})$ over its own averaged
part, and nan with fewer than two values (\code{segment\_mean}):
\[
T=\frac{\sum_k \tfrac12(\kappa_k+\kappa_{k+1})(s_{k+1}-s_k)}{s_{\text{last}}-s_{\text{first}}}.
\]
This gives $T_{\text{tracer}}$ over $[t_0+5,t_1-5]$ and $T_{\text{Medis}}$ over $[m_0+5,m_1-5]$. Because
the profile is computed on the whole line, the values inside the averaged part do not depend on where the
other line ends.

\section{Bifurcation angle, per participant}

\textbf{Step 12: angle in each tree.} One function is used for the cohort feature and for this comparison
(\code{bifurcation\_angle.angle\_with\_reason}). The angle is nan with a reason when the tree has no point
labelled \code{lm}, or when the LAD or the LCx is missing. Otherwise the LAD and LCx from Step 5 must start
at the same point; if they do not, it is nan. On each daughter the points at arc lengths
$5, 5.5,\dots,10$ mm from O are interpolated, and the daughter's direction is the first principal axis of
these 11 points (SVD), pointing from the 5 mm point towards the 10 mm point. The angle is
$\operatorname{atan2}(\lVert u\times v\rVert,\ u\cdot v)$ in degrees. A daughter shorter than 10 mm gives
nan.

\textbf{Step 13: which angles are compared.} A participant's angle pair is compared only when two things
hold. First, the LAD pair and the LCx pair are both aligned in Step 8 (\code{aligned\_whole}); the
averaged-part minimum of Step 9 does not apply here. Second, both daughters pass the near-O check
(\code{near\_o}), which runs only when the first condition holds. For one daughter, each tree's first
10 mm from its own O is resampled every 0.25 mm. The distances run from it to the other tree's left main
(ostium to O, \code{split\_trunks.left\_main}) together with that tree's whole daughter, also resampled, in
both directions. The larger of the two 90th percentiles must be at most 2 mm. The left main is in the
reference because the two O need not coincide. If one tree places O a few mm further along the same vessel,
the start of its daughter lies beside the other tree's left main, and the pair still passes. A daughter
that enters a different vessel within 10 mm still fails. A daughter shorter than 10 mm in either tree gives
nan, which fails the check. A computation that raises is recorded as \code{error:<type>} and counted in
the report.

The distance $o=\lVert O_{\text{tracer}}-O_{\text{Medis}}\rVert$ (the first points of the two LADs) is
saved and reported, not used as a cut.

\section{Daughter asymmetry, per participant}

These rules were fixed on 19 November 2026, before any asymmetry result was seen.

\textbf{Step 14: which participants.} The asymmetry is computed only for participants whose angle is
compared (both conditions of Step 13), so both trees describe the same two daughters. The participant must
also have Medis baseline lumen meshes on the left tree. All of that participant's left-tree lumen files are
read and merged into one mesh. The tracer side is the predicted \code{lumen\_surface\_ostium\_pruned.vtk}.
It is a triangle mesh in the same frame as the centerlines, and it goes with the same centerline files
(\code{lumen\_probe.md}).

\textbf{Step 15: cross-sections.} On each tree's own LAD and LCx, take the 11 depths 5, 5.5, \dots, 10 mm
from that tree's own O. At each depth, the tree's own lumen surface is cut with the plane through the
centerline point $p$. The plane's normal is the direction from the point 1 mm before $p$ to the point 1 mm
after it. Only triangles with a vertex within 10 mm of $p$ take part. The cut gives line pieces, which are
chained into loops, and chains that do not close are dropped. The cross-section is the closed loop that
encloses $p$ in the plane (the smallest, if several do). Its diameter is the area-equivalent
$D=2\sqrt{A/\pi}$. A depth has no diameter in three cases: no closed loop encloses $p$ (the surface ends
there), the cut leaves the 10 mm radius, or the loop also encloses the point where the other daughter's
centerline crosses the plane (both daughters in one loop near O)
(\code{daughter\_ratio.surface\_diameters}).

\textbf{Step 16: strict coverage and the feature.} A daughter's diameter is the median of its 11 values,
and only when all 11 depths have a diameter, on both trees. This matters most for Medis, whose lumen
surface ends where Medis stopped contouring. The asymmetry is $a=\ln(D_{\text{LCx}}/D_{\text{LAD}})$ on
each tree. A participant is compared when all four daughters (LAD and LCx, tracer and Medis) are complete.
Otherwise the reason is recorded: Medis lumen not covering 5--10 mm (checked first), a tracer cut missing,
or \code{error:<type>}.

Medis contoured almost only diseased segments. Every participant compared therefore has plaque at the
bifurcation, where the lumen is narrowed. The plaque-free cohort may agree better, and this comparison
cannot show it.

\section{Agreement}

\textbf{Step 17: differences.} For each feature, the difference is tracer $-$ Medis in the feature's own
unit (1/mm, degrees), a plain difference with no ratio or log.

\textbf{Step 18: statistics.} The bias is the mean difference, $\mathit{SD}$ is the standard deviation
of the differences ($n-1$), and the 95\% limits of agreement are bias $\pm 1.96\,\mathit{SD}$. The 95\%
intervals come from 2000 bootstrap samples of participants: participants are drawn with replacement, all
their pairs come along, and the 2.5th and 97.5th percentiles of the bias and of both limits are taken.
Arteries of one participant are therefore not treated as independent.

\textbf{Step 19: rows reported.} For curvature: the LAD, the LCx, the RCA and all arteries, from the
``ok'' pairs, plus one sensitivity row that also holds the ``offset $>$ 2 mm'' pairs. Those pairs meet the
same length rules as the main rows. For the angle: one row, from the participants of Step 13 with an angle
in both trees. For the asymmetry: the signed
$a$, and the LAD and the LCx diameters in mm as secondary rows, from the participants of Step 16, with the
number of participants left out for each reason.

Next to the limits the report gives the flow (participants and pairs at each step) and the pairs per
artery and reason, as a share of that artery. It also gives the median alignment distances of the pairs
kept, the number of short shared stretches whose two lengths differ more than twofold, the flow of the
angle (aligned on the whole artery, also near O, angle in both, reasons for a missing angle), and the
median and 90th percentile of $o$ with the share above 5 mm. The limits describe only the pairs kept; the
plaque-free cohort uses every trace.

\textbf{Disclosure.} Counts of 1 to 4 are shown as $<5$. A statistic needs at least 5 pairs, a median at
least 10 values and a 90th percentile at least 20. Per-pair rows (with study ids) and the Bland--Altman
plots stay on the server.

\section{What changed on 19 November 2026}\label{sec:changes}

These changes come from the code review, not from the numbers, and they were made after the
18 November results had been seen.
\begin{itemize}
\item Step 5: an \code{lm} label past O no longer moves a daughter's start when the two paths share one.
  Before, the angle could become nan without a warning.
\item Step 8: both directions are measured between the evenly resampled shared stretches. Before, the
  query was the raw points and the reference was the whole other line.
\item Step 9: the averaged-part reason comes before the offset. Before, offset failures skipped it and
  entered the sensitivity row with averages over a few mm.
\item Step 13: the near-O check was added. Its first version (never run on the data) compared with the
  other daughter's first 15 mm only. An O placed more than about 3 mm further along the vessel then failed
  the check by itself, which made it a cut on $o$. The left main is now part of the reference.
\item Steps 14--16: the daughter asymmetry was added (the tracer's lumen was delivered on that day).
\item Step 6 onward: a pair that raises is counted (\code{error:<type>}) instead of stopping the run.
\end{itemize}

\end{document}
